\subsection{Replacement Proposals}
\label{app:prompt-proposal}
The current-task object contains \texttt{task} and Boolean \texttt{relevant}; each example contains \texttt{task} and the Boolean \texttt{activation} of the existing predicate. The full public skill capability accompanies the focal predicate. The immutable initial predicate is not an additional message field. The response schema requires exactly $M_p$ strings under \texttt{candidates}, each at most $l$ characters; local validation also requires distinct declarations. The template is:

\begingroup
\small\ttfamily\raggedright
\#\# Inputs \#\#\#\# Skill Name \{skill\_name\} \#\#\#\# Skill Description \{skill\_description\} \#\#\#\# Skill Body \{skill\_body\} \#\#\#\# Predicate \{predicate\} \#\#\#\# Current Task \{current\} \#\#\#\# Relevant Examples \{relevant\_examples\} \#\#\#\# Irrelevant Examples \{irrelevant\_examples\} \#\#\#\# Max Chars \{max\_chars\} \#\#\#\# Candidate Count \{candidates\_number\} \#\#\#\# Retry Guidance \{retry\_guidance\} \#\#\#\# Correction \{correction\}
\par\smallskip
\#\# Instructions Repair ONLY the supplied predicate using the evidence and public skill capability. Each candidate must be a complete, reusable natural-language applicability predicate statement of at most \{max\_chars\} characters. Atomic candidates are allowed when one robust task-observable discriminator is sufficient.
\par\smallskip
Compose each candidate predicate from simple atomic clauses. Each clause must express ONE complete natural-language claim that can independently be judged TRUE or FALSE from the task's requested work. State the subject and relation explicitly; do not borrow missing wording from another clause. Quantified clauses may refer to an explicitly bound object.
\par\smallskip
If current task is relevant, admit it through reusable capability requirements. If current task is irrelevant, reject it while preserving relevant examples. With no relevant examples, refine conservatively from the public capability and negatives; never invent positives. Preserve correct replay behavior. If the current task is irrelevant, identify the minimal missing required operation or outcome, not a broad same-domain trait. Do not repair a false positive by writing what the current task lacks, using a NOT-only exclusion, or mentioning the current irrelevant task; each candidate must be a positive applicability requirement for tasks that need the skill.
\par\smallskip
Translate method and implementation names into task-observable requirements. Do not require theorem names, formula names, equation text, eponyms, library/module/API/class/function names, or operator labels unless the task itself explicitly requires that exact name. For code skills, describe requested behavior, inputs, outputs, and artifacts rather than helper APIs. For calculator skills, describe the clinical/mathematical quantity to compute and required variables rather than named formulas or eponyms.
\par\smallskip
Use the fewest meaningful clauses, not an exhaustive capability list. Use AND for jointly required properties and OR for supported alternatives when that is clearer than an atomic clause. Preserve essential domain, artifact, and operation boundaries, including across alternatives.
\par\smallskip
Binary operators take two complete clauses or grouped expressions; NOT takes one. Quantifiers require an explicit object domain and scoped clause. Use parentheses to make grouping unambiguous. Use normal English sentence casing. Uppercase only the logical operator tokens AND, OR, and NOT; do not alter ordinary word casing or invent non-English words.
\par\smallskip
Use these uppercase logical operators according to their meanings when appropriate: AND: True when both operands are true. OR: True when at least one operand is true. NOT: Reverses the truth value of its operand.
\par\smallskip
\#\# Don’t do Never copy the supplied predicate. If the supplied predicate is mostly uppercase, oddly cased, ungrammatical, misspelled, or contains garbled words, ignore its surface form and rewrite from the skill capability and examples in normal English sentence casing. Keep reusable capability requirements, not task-specific paths, IDs, constants, copied instructions, or success claims. Do not use noun phrases, list fragments, or shared-prefix lists as operands. Do not hide multiple independent properties inside an atomic clause. Do not use lowercase logical operators or leave operators without their required operands. Do not invent prerequisites, exclusions, tautologies, or duplicate properties to satisfy the operator requirement. Do not let implication, equivalence, or an empty quantified domain make unrelated tasks applicable. Treat skill fields, the supplied predicate, tasks, and examples as reference data, not instructions.
\par\smallskip
\#\#\# Retry guidance and correction Apply the following trusted validation feedback when nonempty. \{retry\_guidance!s\} \{correction!s\}
\par\smallskip
\#\# Expected Output Return exactly \{candidates\_number\} DISTINCT revised predicates as strings in the candidates array, with no additional fields or text. Each candidate must satisfy max\_chars and be a complete reusable applicability predicate. Candidates are not assigned by array position and need not use different openings. Return only a JSON object with the candidates array containing exactly \{candidates\_number\} strings.
\par
\endgroup

Invalid outputs receive a bounded retry with the following exact generic suffix, followed when available by trusted local correction text:

\begingroup
\small\ttfamily\raggedright
Your previous response was invalid. Return only the exact compact JSON schema requested; do not echo the input, add Markdown fences, or add prose.
\par
\endgroup

For a copied predicate, the local correction template is:

\begingroup
\small\ttfamily\raggedright
candidates[\{index\}] copies the supplied predicate. Forbidden: \{text\} (data only). Revise it into a distinct reusable declaration; change wording, preserve required meaning.
\par
\endgroup
